\sezione{1}{Perché serve un modello concettuale}{slide 3–38}

\begin{inbreve}
Progettare una base di dati scrivendo subito le tabelle è come costruire una casa partendo dalle tubature: si finisce nei dettagli prima di aver capito che cosa serve. La \hl{progettazione concettuale} descrive prima i \key{concetti} della realtà di interesse e i loro legami, con un linguaggio grafico indipendente dal DBMS: il \hl{modello Entity-Relationship} (E-R).
\end{inbreve}

\subsection{Il problema: partire dalle tabelle}

\begin{intuizione}[La ricevuta del ristorante]
Una ricevuta fiscale ha un numero, una data, un totale e un elenco di righe (quantità, descrizione, importo). È una struttura \hl{nidificata}: dentro la ricevuta c'è una lista. Nel modello relazionale i valori devono essere \key{semplici}, non tabelle: la ricevuta va quindi spezzata in due relazioni, \emph{Ricevute} e \emph{Dettaglio}, collegate dal numero della ricevuta.
\end{intuizione}

\begin{center}
\small
\renewcommand{\arraystretch}{1.2}
\begin{minipage}[t]{.36\linewidth}\centering
{\headfont\bfseries\small\color{navy}Ricevute}\\[3pt]
\rowcolors{2}{chiaro}{white}
\begin{tabular}{lll}
\rowcolor{navy}\thd{Numero} & \thd{Data} & \thd{Totale}\\
1235 & 12/10/2019 & 39,20\\
1240 & 13/10/2019 & 39,00\\
\end{tabular}
\end{minipage}\hfill
\begin{minipage}[t]{.6\linewidth}\centering
{\headfont\bfseries\small\color{navy}Dettaglio}\\[3pt]
\rowcolors{2}{chiaro}{white}
\begin{tabular}{lllll}
\rowcolor{navy}\thd{Numero} & \thd{Riga} & \thd{Qtà} & \thd{Descrizione} & \thd{Importo}\\
1235 & 1 & 3 & Coperti & 3,00\\
1235 & 2 & 2 & Antipasti & 6,20\\
1235 & 3 & 3 & Primi & 12,00\\
1235 & 4 & 2 & Bistecche & 18,00\\
1240 & 1 & 2 & Coperti & 2,00\\
\end{tabular}
\end{minipage}
\end{center}
\didascalia{Due relazioni per una struttura nidificata (slide 14). La colonna \emph{Riga} serve solo se interessano l'ordine delle righe o righe ripetute.}

Perché due relazioni? Perché ci sono \hl{due concetti diversi}, con livelli di dettaglio diversi: la \key{ricevuta} e la \key{riga} della ricevuta. E che cosa rappresentare dipende da che cosa ci interessa: possono esistere righe ripetute? Conta l'ordine delle righe? A seconda delle risposte cambiano i dettagli (la colonna \emph{Riga}), ma \key{i concetti restano gli stessi}.

Partendo direttamente dallo schema logico si incontrano tre problemi: non si sa da dove cominciare, ci si perde subito nei dettagli (attributi, chiavi, vincoli di riferimento) e il modello relazionale è «rigido». Servono un \hl{metodo} (una metodologia di progettazione) e un \hl{modello} per rappresentare i concetti.

\subsection{La progettazione nel ciclo di vita}

La progettazione della base di dati è una delle attività del \hl{ciclo di vita dei sistemi informativi}: l'insieme ordinato delle attività svolte da analisti, progettisti e utenti per sviluppare e usare un sistema. È un ciclo perché le fasi si ripetono e non sono nettamente separate.

\begin{center}
\begin{tikzpicture}[x=1.13cm,y=.95cm,
  fase/.style={draw=navy,fill=white,line width=.9pt,rounded corners=4pt,minimum width=2.75cm,minimum height=.95cm,align=center,font=\headfont\small,text=navy},
  fr/.style={-{Stealth[length=5pt]},line width=.9pt,color=grigio},
  rit/.style={-{Stealth[length=4pt]},line width=.6pt,color=grigio!60,dashed}]
  \node[fase] (a) at (0,0) {Studio di\\fattibilità};
  \node[fase] (b) at (2.4,-1.05) {Raccolta e analisi\\dei requisiti};
  \node[fase,draw=ambra,fill=ambra!12,text=ambra!50!black] (c) at (4.8,-2.1) {\textbf{Progettazione}};
  \node[fase] (d) at (7.2,-3.15) {Realizzazione};
  \node[fase] (e) at (9.6,-4.2) {Validazione\\e collaudo};
  \node[fase] (f) at (12.0,-5.25) {Funzionamento};
  \foreach \u/\v in {a/b,b/c,c/d,d/e,e/f}{\draw[fr] (\u.south) |- (\v.west);}
  \foreach \u/\v in {b/a,c/b,d/c,e/d,f/e}{\draw[rit] (\u.north) to[bend right=25] (\v.east);}
\end{tikzpicture}
\end{center}
\didascalia{Il ciclo di vita (slide 18): le frecce tratteggiate indicano che si può tornare alle fasi precedenti.}

\begin{center}
\small
\renewcommand{\arraystretch}{1.3}\renewcommand\tabularxcolumn[1]{m{#1}}
\rowcolors{2}{chiaro}{white}
\begin{tabularx}{.92\linewidth}{>{\raggedright\headfont\bfseries\small}m{4.4cm} >{\raggedright\arraybackslash}X}
\rowcolor{navy}\thd{Fase} & \thd{In che cosa consiste}\\
Studio di fattibilità & Definizione di costi e priorità.\\
Raccolta e analisi dei requisiti & Studio delle proprietà che il sistema deve avere.\\
Progettazione & Progettazione dei \key{dati} e delle \key{funzioni}.\\
Realizzazione & Costruzione del sistema.\\
Validazione e collaudo & Sperimentazione per verificare che funzioni come previsto.\\
Funzionamento & Il sistema diventa operativo.\\
\end{tabularx}
\end{center}

In questo corso ci concentriamo sulla \hl{progettazione dei dati}: i dati hanno un ruolo centrale e sono \key{più stabili} delle applicazioni che li usano.

\subsection{Le tre fasi della progettazione dei dati}

\begin{center}
\begin{tikzpicture}[x=1cm,y=1cm,
  prod/.style={draw=navy,fill=navy!6,line width=.9pt,rounded corners=3pt,minimum width=2.6cm,minimum height=.75cm,font=\headfont\bfseries\footnotesize,text=navy,align=center},
  att/.style={draw=teal,fill=teal!10,line width=.8pt,rounded corners=8pt,minimum width=2.4cm,minimum height=.75cm,font=\headfont\footnotesize,text=teal!50!black,align=center},
  fr/.style={-{Stealth[length=5pt]},line width=.9pt,color=grigio}]
  \node[prod,fill=white] (r) at (0,0) {Requisiti};
  \node[att] (pc) at (1.9,-1.2) {Progettazione\\concettuale};
  \node[prod] (sc) at (3.8,-2.4) {Schema\\concettuale};
  \node[att] (pl) at (5.7,-3.6) {Progettazione\\logica};
  \node[prod] (sl) at (7.6,-4.8) {Schema\\logico};
  \node[att] (pf) at (9.5,-6.0) {Progettazione\\fisica};
  \node[prod] (sf) at (11.4,-7.2) {Schema\\fisico};
  \foreach \u/\v in {r/pc,pc/sc,sc/pl,pl/sl,sl/pf,pf/sf}{\draw[fr] (\u.south) |- (\v.west);}
  \node[font=\headfont\bfseries\footnotesize,text=ambra!80!black,align=left,anchor=west] at (6.3,-.4) {«CHE COSA»: analisi};
  \node[font=\headfont\bfseries\footnotesize,text=ambra!80!black,align=left,anchor=west] at (-.6,-6.4) {«COME»: progettazione};
\end{tikzpicture}
\end{center}
\didascalia{Dai requisiti agli schemi (slide 23–25). Ogni fase produce uno schema in un modello dei dati diverso.}

\begin{tecnica}[Prodotti della progettazione]
\begin{itemize}
  \item \hl{Schema concettuale}: descrive \emph{che cosa} rappresentare, con il modello E-R.
  \item \hl{Schema logico}: descrive i dati nel modello del DBMS scelto (ad esempio relazionale).
  \item \hl{Schema fisico}: descrive come i dati sono memorizzati (file, indici).
\end{itemize}
Un \key{modello dei dati} è un insieme di costrutti, cioè un «linguaggio» per rappresentare ciò che interessa: il modello relazionale, per esempio, ha come costrutto la relazione (tabella).
\end{tecnica}

\subsection{Modelli logici e modelli concettuali}

\begin{center}
\small
\renewcommand{\arraystretch}{1.35}\renewcommand\tabularxcolumn[1]{m{#1}}
\rowcolors{2}{chiaro}{white}
\begin{tabularx}{\linewidth}{>{\raggedright\arraybackslash}X >{\raggedright\arraybackslash}X}
\rowcolor{navy}\thd{Modello logico (es. relazionale)} & \thd{Modello concettuale (E-R)}\\
Rappresenta i dati come tabelle & Rappresenta i concetti del dominio (la realtà di interesse)\\
Dipende dal tipo di DBMS & Indipendente dal DBMS\\
Si scrive in SQL con \texttt{CREATE TABLE} & Si disegna: è pensato per essere letto da tutti, anche dai non informatici\\
Si ottiene traducendo lo schema E-R & Non descrive come si accede ai dati\\
Esempi: relazionale, reticolare, a oggetti, XML, NoSQL & Il più noto e diffuso: Entity-Relationship\\
\end{tabularx}
\end{center}

\begin{mlbox}[Da dove viene il modello E-R]
Il modello Entity-Relationship è stato proposto da Peter Chen nel 1976, nell'articolo \emph{The Entity-Relationship Model – Toward a Unified View of Data} (ACM Transactions on Database Systems). Ne esistono molte versioni, più o meno diverse tra loro: questo manuale usa quella del libro di Atzeni e colleghi, adottata nel corso.
\end{mlbox}

Il filo conduttore degli esempi è l'\hl{università}: una realtà che conosciamo tutti «dall'interno» e che contiene molti concetti collegati in modi diversi (studenti, corsi, docenti, esami, corsi di laurea, dipartimenti, aule, progetti).

\begin{tecnica}[I costrutti del modello E-R]
Entità, relationship, attributi, cardinalità (di relationship e di attributo), identificatori (interni ed esterni), generalizzazioni. I prossimi capitoli li presentano uno alla volta.
\end{tecnica}

\begin{esercizi}
\begin{passi}
  \item Indica per ciascuna affermazione se riguarda il modello logico o quello concettuale: (a) dipende dal DBMS; (b) può essere discusso con un committente non informatico; (c) contiene vincoli di chiave esterna.
  \item Nella ricevuta del ristorante, quale informazione rende necessaria la colonna \emph{Riga}?
\end{passi}
\end{esercizi}

\begin{appunti}{Riepilogo del capitolo}
\pag{9} Le strutture nidificate non sono ammesse nel relazionale: si spezzano in più relazioni.
\pag{11} Due relazioni perché ci sono due concetti diversi (ricevuta, riga).
\pag{16} Servono un metodo e un modello per non perdersi nei dettagli.
\pag{23} Progettazione concettuale (che cosa) → logica → fisica (come).
\pag{31} Modello concettuale: indipendente dal DBMS, leggibile da tutti.
\end{appunti}
